% Companion note 6 to "Relational Progression as an Operational Foundation for GR".
% Prepared with Claude (Anthropic) at the author's request.
\documentclass[11pt]{article}
\usepackage[margin=1in]{geometry}
\usepackage{amsmath,amssymb}
\usepackage{booktabs,array}
\usepackage{xcolor}
\usepackage[colorlinks=true,linkcolor=blue!50!black,citecolor=blue!50!black,urlcolor=blue!50!black]{hyperref}
\newcommand{\dd}{\mathrm{d}}
\newcommand{\st}[1]{\textsf{\small[#1]}}
\title{Entanglement Equilibrium for a Fermion Environment:\\ the Induced-Gravity Lattice Test Repeated for Dirac Fields\\[4pt]
\large Companion note 6 to the consolidated draft}
\author{Scott A. Jordan}
\date{Draft, 24 November 2026}
\begin{document}
\sloppy
\maketitle
\begin{abstract}
Companion note~5 found that a free scalar cannot both induce $G$ and satisfy entanglement equilibrium. At minimal coupling it gives $G=1/(4\eta)$, but its ball entropy carries a term $\propto R^2\delta\langle\phi^2\rangle$ that the geometry cannot cancel. At conformal coupling equilibrium is restored, but the induced $G$ vanishes. Here we repeat the test for a free massless Dirac fermion in the same induced-gravity setup: the ruler $f(r)$ and the lapse (progression factor) are classical fields whose only dynamics is the fermion's own vacuum Killing energy, with the regulator-independent part isolated by the same Pauli--Villars combination. Lattice fermion doubling is handled with a staggered radial lattice and a covariant radial derivative; a naive lattice is run alongside and doubles, as it should. With the conditions written down before running, all pre-registered tests pass.
(i) The fermion's ball modular Hamiltonian is the CHM operator $2\pi\int w\,T_{00}$ built from the symmetric (gravitating) stress tensor, to about $1\%$. It has no extra bulk term (excluded exactly by chiral symmetry) and no edge term.
(ii) The fermion induces a nonzero $G$, and $G_{\rm ind}=1/(4\eta_D)$ with the fermion's own entanglement coefficient to within $3\%$ at two mass scales.
(iii) The induced action has the Einstein form, with the lapse entering linearly.
So, for the fermion in this model, entanglement equilibrium holds for all states and is derived. $G$ is fixed by the environment in its regulator-independent part, but not in its cutoff-scale part.
\end{abstract}

\section{Why fermions, and what changes}
In note~5 the scalar's obstruction came from $\phi^2$, a scalar operator of dimension $\Delta=1<d/2$, whose ball term scales as $R^{2\Delta}=R^2$ and dominates the $R^4$ energy term. A Dirac fermion has no gauge-invariant scalar of dimension below $3$ ($\bar\psi\psi$ has $\Delta=3$, giving $R^6$), it admits no nonminimal curvature coupling of dimension $\le4$, and it has no contact term in its entanglement entropy~\cite{Kabat}. The covariant (heat-kernel) expectations for one Dirac field, relative to one minimally coupled real scalar, are $\Lambda_D=-4\Lambda_S$, $C_D=2C_S$, $\eta_D=2\eta_S$. In the Pauli--Villars scheme of note~5, with $\mathcal A=6\ln2-3\ln3$, this means
\begin{equation}
\eta_D=\frac{\mathcal A m^2}{24\pi},\qquad C_D=-\frac{1}{16\pi G_{\rm ind}}=-\frac{\mathcal A m^2}{96\pi^2},\qquad\text{so}\quad -4\pi C_D=\eta_D .
\end{equation}

\section{The lattice Dirac field, and doubling}
Separating the angular dependence (quantum number $\kappa=\pm1,\pm2,\dots$, degeneracy $2|\kappa|$) leaves a two-component radial problem $H_\kappa=\begin{pmatrix}m&B^\dagger\\B&-m\end{pmatrix}$ with $B=\partial_r+\kappa/f(r)$ on $\dd\ell^2=\dd r^2+f^2\dd\Omega^2$ (the spin connection is absorbed by $\psi=\chi/f$). The lapse enters as $H_N=\sqrt N\,H\sqrt N$, which reproduces the $\tfrac12\alpha\cdot\nabla N$ spin-connection term.

\textbf{Doubling.} A naive discretisation (both components on the same sites, central differences) has a second massless branch at lattice momentum $\pi$. We instead (a) stagger: the upper component lives on sites $r=j$ and the lower on $r=j+\tfrac12$, so the derivative is a nearest-neighbour difference with a single zero at $k=0$; and (b) discretise $B=e^{-\Phi}\partial_re^{\Phi}$ with $\Phi'=\kappa/f$ exactly, $(Bu)_{j+\frac12}=e^{-\Phi_{j+\frac12}}(e^{\Phi_{j+1}}u_{j+1}-e^{\Phi_j}u_j)$. Averaging $\kappa/f$ instead would leave a low-energy mode at $k=\pi$ for large $\kappa/r$; with (b) the gap is $\ge2\sinh(\kappa/2f)$ at every lattice momentum. (c) The chain is terminated so that the regular solution is kept. A first version with the same termination for both signs of $\kappa$ had a spurious near-zero mode $u\propto r^{-\kappa}$ localized at the origin, which made the entropy erratic; it was found and removed. Checks:
\begin{center}\small\begin{tabular}{lccc}\toprule
& staggered, covariant & naive & continuum\\\midrule
positive levels below $E=0.5$ ($\kappa=-1$, $L=60$) & 9 & 19 & 9.6\\
vacuum energy density $\Lambda_D/\Lambda_S$ (PV, $m=0.2$) & $-3.88$ & $-8.15$ & $-4$ (doubler: $-8$)\\\bottomrule\end{tabular}\end{center}
The staggered field is a single Dirac fermion. Its $3\%$ deficit in $\Lambda$ is the $O((Ma)^2)$ lattice artifact; the scalar's was $+1.2\%$.

\section{Pre-registered tests}
Written before any fermion run (file \texttt{PREREG\_note6.md}):
\textbf{P0} no doubling ($\Lambda_D/\Lambda_S\approx-4$, not $-8$).
\textbf{P1F} induced action of the Einstein form; lapse--lapse term at $O(k^2)$ below $10\%$ of the Einstein--Hilbert scale and falling with width.
\textbf{P2F} $-4\pi C_D/\eta_D=1$ within $10\%$, failing otherwise.
\textbf{P4F} ball modular Hamiltonian $=2\pi\cdot\tfrac12\{W,H\}$ (symmetric $T_{00}$), with any extra $\psi^\dagger\psi$ or $\bar\psi\psi$ bulk term below $0.1$ of the scalar's $\tfrac{\pi}{R}\phi^2$ term, failing otherwise.

\section{Results}
\paragraph{P4F: the ball modular Hamiltonian} (massless, 100-digit precision, $n=24$, $R=24.75$; 30 smooth probes inside the ball plus 6 reaching the sphere). Median $Q_{\rm exact}/Q_{\rm CHM}$:
\begin{center}\small\begin{tabular}{lcccc}\toprule
sector $\kappa$ & $-1$ & $+1$ & $-2$ & $+2$\\\midrule
$N=96$: interior / edge probes & 1.029 / 1.008 & 1.029 / 1.008 & 1.014 / 1.003 & 1.007 / 1.003\\
$N=144$: interior / edge probes & 1.011 / 1.003 & 1.011 / 1.003 & 1.004 / 1.001 & 1.002 / 1.001\\\bottomrule\end{tabular}\end{center}
The residual falls as the box grows, as the scalar's finite-box factor did. Probes that are large at the entangling sphere agree to within $0.1$--$0.8\%$, so there is \emph{no edge term} (contrast note~5, where the scalar's edge term was decisively absent from a form that would have required it). Probes living on one component only ($\psi^\dagger\psi$- or $\bar\psi\psi$-type) give $Q_{\rm exact}\approx10^{-28}$ against an expected scale of order $1$. The massless lattice fermion has a sublattice (chiral) symmetry that makes the reduced correlation matrix $\tfrac12$ on the diagonal blocks, so these terms are forbidden \emph{exactly}, not merely small. In the scalar the corresponding term was $4$--$77\%$ of the total. \textbf{Pass.} (In $\kappa=+2$ a few probes whose CHM value nearly cancels scatter by up to $\pm28\%$; the median and the edge probes are unaffected.)

\paragraph{Entanglement coefficient} (flat lattice, $R=20.75$--$50.75$, fit $S=\eta\,4\pi R^2+b\ln R+c$):
$\eta_D/{\rm prediction}=1.006$ ($m=0.14$), $1.014$ ($m=0.20$), $1.027$ ($m=0.28$); $\eta_D/\eta_S=2.04$ (prediction $2$). The drift with $m$ is the $(Ma)^2$ artifact.

\paragraph{P2F, P1F: induced $G$.} Geometry runs: bump at $r_0=150$, wall at $360$, widths $w=12$--$40$, $\epsilon=0.02$. The second-order $k^0$ coefficient is fitted as a free constant, as in note~5. Its offset from the first-order $\Lambda$ ($2$--$5\times10^{-4}$) changes when the angular truncation is changed, so it is a truncation artifact rather than physics.
\begin{center}\small\begin{tabular}{lccc}\toprule
& $m=0.20$ & $m=0.28$ & $m=0.20$, truncation $l/r>10$\\\midrule
$C_D/C_{\rm pred}$ (four fit forms) & 1.007--1.029 & 0.994--1.011 & 0.967--1.031\\
$-4\pi C_D/\eta_D$ (P2F) & 0.994--1.015 & 0.968--0.984 & 0.954--1.017\\\bottomrule\end{tabular}\end{center}
Unlike the minimally coupled scalar, whose direct curvature coefficient had not converged ($10$--$30\%$ high), the fermion converges at the widths used. A massless Dirac field is conformal, so its four-derivative action contains no $R^2$ term (only Weyl$^2$ and the Euler density), which is why the contamination is small. \textbf{P2F passes}: $G_{\rm ind}=1/(4\eta_D)$ to within $3\%$.

Lapse (progression factor), $m=0.2$:
\begin{center}\small\begin{tabular}{lcccc}\toprule
$w$ & 16 & 24 & 32 & 40\\\midrule
lapse--lapse / EH scale & $-0.055$ & $-0.013$ & $+0.006$ & $+0.016$\\
lapse--ruler cross term, $C_\times/C_{\rm pred}$ & 0.934 & 0.966 & 0.981 & 0.993\\\bottomrule\end{tabular}\end{center}
There is no $(\nabla N)^2$ term at the few-percent level, and the clock--ruler cross term converges to the same $G$. \textbf{P1F passes.}

\emph{Precision.} With $\epsilon=0.02$, the lapse--lapse term first showed a spurious growth at high angular momentum. This was rounding noise from the large staggered couplings ($\le e^{5}$), amplified by $1/\epsilon^2$. It disappears for $\epsilon=0.05$ or a tighter truncation ($l/r>5$); with that truncation the partial sums converge by $l\approx400$--$1000$ to within $4\times10^{-4}$. All final geometry runs use the $l/r>5$ truncation (the entropy and $\Lambda$ runs, which involve no $1/\epsilon^2$ amplification, used $l/r>10$). Results are unchanged in $\epsilon$ and in angular cutoff. Finite box: the modular-Hamiltonian residual falls from $3\%$ to $1\%$ between $N/R=4$ and $6$.

\section{Status}
For a free Dirac fermion in the induced-gravity model:
\begin{itemize}
\item the geometry's dynamics, generated by the fermion's vacuum energy, is Einstein-like, with $G_{\rm ind}=1/(4\eta_D)$ (mass-dependent part);
\item the fermion's ball entropy responds to its gravitating (symmetric) stress tensor with the CHM weight, with no $R^2$ obstruction and no edge term.
\end{itemize}
Jacobson's argument therefore closes for this environment: for every first-order state change, the entropy change is balanced by the area change that the induced dynamics produces. \textbf{Entanglement equilibrium: derived (numerical) for a fermion environment in induced gravity, regulator-independent part.} The scalar result of note~5 stands: a scalar environment violates it at minimal coupling. Equilibrium is therefore a property of the field content, not a universal principle. \textbf{$G$:} fixed by the environment in its regulator-independent (mass-dependent) part, $1/(4G)=\eta$. The cutoff-scale part ($\propto a^{-2}$) is still not computed, so the absolute value of $G$ is not predicted. Status: Conditional.

\paragraph{Limits.} Static geometries; classical (mean-field) geometry; free fields; regulator-independent part only; radial lattice with continuum angular directions; Pauli--Villars masses up to $Ma\approx0.5$ (artifacts $1$--$3\%$); box of radius $360$ with bumps at $r_0=150$, widths up to $40$.

\paragraph{Next.} (i) Mixed environments: a scalar plus fermions, to see whether the scalar's $R^2$ obstruction can be cancelled (it cannot by fermions, which have no $\Delta=1$ operator, so a realistic environment must avoid light minimally coupled scalars or supply a compensating mechanism). (ii) Gauge fields, which have a known contact-term subtlety. (iii) A covariant regulator with a statistical entropy (for example a random lattice), so that the cutoff-scale part of $1/G$ can be compared with $\eta$.

\section*{Code}
\texttt{fermion\_common.py} (staggered covariant lattice Dirac field), \texttt{fermion\_doubling.py} (P0, naive comparison), \texttt{fermion\_modular.py} (P4F, \texttt{mpmath}), \texttt{fermion\_eta.py}, \texttt{fermion\_G.py} (geometry side; args $m$, $0$, mode, widths), \texttt{fnncheck2.py} (precision check), \texttt{analyse\_FG.py}, \texttt{analyse\_feta.py}, \texttt{PREREG\_note6.md}; raw results in \texttt{note6\_results/}.

\begin{thebibliography}{9}
\bibitem{Kabat} D.~Kabat, ``Black hole entropy and entropy of entanglement,'' \emph{Nucl.\ Phys.\ B} \textbf{453}, 281 (1995).
\bibitem{J2016} T.~Jacobson, ``Entanglement equilibrium and the Einstein equation,'' \emph{Phys.\ Rev.\ Lett.} \textbf{116}, 201101 (2016).
\bibitem{CGM2016} H.~Casini, D.~A.~Galante, R.~C.~Myers, \emph{JHEP} \textbf{03}, 194 (2016), arXiv:1601.00528.
\end{thebibliography}
\end{document}
